\documentclass[11pt]{tex/lawoutline}
\usepackage{amssymb}
\usepackage{enumitem}

\course{Legislation and Regulation}
\student{Prof: R. Doerfler}
\semester{Fall 2026}

\begin{document}

\maketitle
\tableofcontents

\chapter{The Legislative Process and Theories of Interpretation (Legislation)}

\section{Text and Purpose}

\section{Ordinary Meaning, Special Meaning, and
Context}

\section{The Legislative Process and Legislative History}

\chapter{Canons of Construction (Legislation)}

\section{Linguistic Canons}

\section{Substantive Canons}
\definition{Substantive Canons}
{These canons ask interpreters to put a thumb on the scale in favor of some value or policy that courts have identified as worthy of protection.}
\begin{itemize}
    \item In order for these canons to do any work, it must be the case that without the use of the canon, the court would have reached a different conclusion. 
    \item Substantive canons may also be understood as techniques the courts employ in order to advance particular values that the \textit{courts} have decided are important (rather than Congress). 
\end{itemize}
\question{Legitimacy}{Under what conditions is it legitimate for courts to develop and apply substantive canons?}

\subsection{Avoiding Constitutional Questions}

\rulebox{Avoidance Canon}
{
\begin{itemize}
    \item Doctrine that courts should construe statutes to avoid serious constitutional problems. 
    \item Historically, $\exists$ a strong presumption in favor of upholding statutes as constitutional by the Supreme Court. 
    \item \textbf{Classical Avoidance Canon:} Courts are supposed to interpret statutes in order to avoid constructions that would \textit{actually} be unconstitutional. 
    \item Classical avoidance is designed to avoid judicial invalidation of statutes on constitutional grounds, and it necessitates that courts decide the constitutional issue. 
    \item \textbf{Modern Avoidance Canon:} The presence of a serious constitutional doubt or question about one possible construction of a statute is a sufficient reason to adopt a different construction, as long as it is fairly possible. (J. Brandeis)
    \item Modern avoidance allows courts to avoid deciding the constitutional question completely. 
\end{itemize}
}

\casebox{\textit{NLRB v. Catholic Bishop of Chicago}}
{Supreme Court, 1979}
{
\textbf{Facts}
\begin{itemize}
    \item The Catholic Bishop of Chicago operated two Catholic ``minor seminaries,'' and the Diocese of Fort Wayne-South Bend operated five Catholic high schools.
    \item The schools taught both secular subjects, such as mathematics, English, and science, and mandatory religious subjects. 
    \item Unions sought to represent the schools' lay teachers. Religious faculty and various administrative and nonteaching employees were not included in the proposed bargaining units.
    \item The NLRB supervised union elections, but the schools refused to recognize or bargain with the unions, arguing that the NLRB lacked jurisdiction over church-operated schools and that NLRB regulation would violate the First Amendment's Religion Clauses.
    \item The NLRB ordered the schools to recognize and bargain with the unions.
    \item The Seventh Circuit refused to enforce the NLRB's orders, reasoning that the Board's jurisdictional standard was unclear and that applying the NLRA to church-operated schools raised serious First Amendment concerns.
    \item The Supreme Court affirmed the Seventh Circuit. It held that the NLRA should not be interpreted to give the NLRB jurisdiction over teachers in church-operated schools absent a clear expression from Congress.
\end{itemize}

\textbf{Question}
\begin{itemize}
    \item Do teachers in schools operated by a church to teach both religious and secular subject within the jurisdiction granted by the NLRA?
    \item If so, is it a violation of the Religion Clauses of the First Amendment?
    \item Supreme Court ruled \textbf{No} on the first question so that they did not have to answer the latter question. 
\end{itemize}

\textbf{Discussion, Majority (J. Burger); Dissent (J. Brennan)}

\begin{enumerate}

    \item \textbf{Does NLRA's language cover church-operated schools?}

    \begin{itemize}
        \item \textbf{Majority:} The NLRA defines ``employer'' broadly and gives the NLRB broad jurisdiction, but the statutory language does not expressly state whether church-operated schools are covered.
        \item The Court acknowledged that the Act contains broad language, but reasoned that the absence of an explicit reference to church-operated schools was important because applying the Act would raise serious constitutional questions.
        
        \item \textbf{Dissent:} The Act's language is not ambiguous. Section 2(2) defines ``employer'' to include all employers except for specifically listed exceptions.
        \item Church-operated schools are not among those exceptions. Therefore, the dissent argued that the Court was effectively adding a new exception that Congress had not enacted.
        \item The dissent characterized the majority's interpretation as substituting judicial amendment for statutory interpretation.
    \end{itemize}

    \item \textbf{The proper form of constitutional avoidance canon}

    \begin{itemize}
        \rulebox{Constitutional Avoidance}
        {When a statute is reasonably capable of more than one interpretation, courts should adopt the interpretation that avoids serious constitutional questions.}

        \item \textbf{Majority:} Courts should not interpret a statute in a way that creates serious constitutional problems unless Congress has clearly expressed its intention to require that result.
        \item Here, applying the NLRA to church-operated schools would create a significant risk of violating the Religion Clauses because the NLRB would have to regulate relationships closely connected to the schools' religious mission.
        \item The Court therefore required an ``affirmative intention'' from Congress before reading the Act to cover lay teachers in church-operated schools.
        \item Because the Court found no such clear expression, it interpreted the Act not to apply.

        \item \textbf{Dissent:} Justice Brennan agreed that constitutional avoidance is a legitimate and important principle, but argued that the majority used the wrong version of the rule.
        \item They should ask whether a constitutional question can be avoided through a construction that is ``fairly possible'' or reasonable, but the majority instead required a clear affirmative statement from Congress before applying a broadly worded statute to an institution whose coverage might raise constitutional concerns.
        \item The dissent argued that this new formulation gives courts too much power to rewrite statutes whenever they have constitutional doubts.
        \item Congress ordinarily does not list every category of person or institution covered by a broadly worded regulatory statute. Requiring Congress to do so would undermine the statute's general language.
    \end{itemize}

    \item \textbf{Statutory history and congressional intent}

    \begin{itemize}
        \item \textbf{Majority:} The legislative history did not show that Congress specifically considered church-operated schools when it enacted the NLRA.
        \item The original 1935 debates focused primarily on workers in private industry and industrial recovery, not on religious schools or religious employers.
        \item Later amendments did not demonstrate a clear congressional decision to include church-operated schools within the NLRB's jurisdiction.
        \item Although Congress had considered exemptions for certain nonprofit institutions, the majority found no clear statement addressing church-operated schools.
        \item Therefore, the legislative history did not supply the affirmative expression necessary to justify an interpretation that would create serious First Amendment questions.

        \item \textbf{Dissent:} The legislative history supported inclusion rather than exclusion.
        \item In 1947, Congress considered a broad exemption for nonprofit, charitable, religious, and educational organizations but did not enact it.
        \item Congress instead created a narrower exemption for nonprofit hospitals, which showed that Congress knew how to create exemptions when it wanted to do so.
        \item Congress later repealed even the nonprofit-hospital exemption in 1974, which is further evidence that Congress intended the NLRA to reach nonprofit employers unless they were expressly excluded.
        \item Congress also rejected proposals to create special religious exemptions, which shows that Congress deliberately declined to enact the exemption.
    \end{itemize}

    \item \textbf{Would NLRB having jurisdiction be unconstitutional?}

    \begin{itemize}
        \item \textbf{Majority:} Teachers play a uniquely important role in church-operated schools, because their conduct and teaching may also communicate or reinforce the school's religious values.
        \item Mandatory collective bargaining could affect many aspects of school administration, including the terms and conditions of teachers' employment.
        \item NLRB might have to decide whether a school's employment decision was motivated by religious doctrine or by unlawful antiunion animus.
        \item Even the process of investigating and deciding unfair-labor-practice claims could require government officials to inquire into the school's religious mission, which could raise constitutionality problems.

        \item \textbf{Dissent:} Constitutional concern did not justify excluding the schools from the Act without first applying the statute's ordinary meaning.
        \item The Court should not avoid a difficult constitutional question by adopting a statutory interpretation that contradicted the Act's language and history.
        \item Relied on prior cases applying the NLRA to nonprofit organizations, including the Associated Press, even when those organizations raised First Amendment concerns.
        \item These precedents suggested that an employer's nonprofit or religious association did not automatically remove it from the NLRA's coverage.
        \item Acknowledged that the constitutional issue might be difficult, but argued that the Court should still address it directly rather than create a new statutory exception.
    \end{itemize}
\end{enumerate}
}

\subsection{State Sovereignty}

\subsection{Preemption of State Law}

\subsection{Scope and Application of Substantive Canons}

\section{In Class Problems}
\lincanon{Rule of Last Antecedent}
{A limiting phrase ordinarily modifies only the noun or phrase immediately preceding it.}

\lincanon{Series-Qualifier Principle}
{When a modifier follows a series of parallel terms, the modifier ordinarily applies to every term in the series when that reading is natural.}

\subcanon{Rule of Lenity}
{Ambiguity in a criminal statute should generally be resolved in favor of the defendant.}

\casebox{\textit{Lockhart v. United States}}
{Supreme Court, 2016}
{
\textbf{Facts}
\begin{itemize}
    \item Lockhart was convicted of sexual abuse in the 1st degree, which involved his adult girlfriend. 
    \item Separately, Lockhart was indicted (and pleaded guilty) for attempting to receive child pornography that violated 18 U.S.C Section 2252(a)(2) and for possession of child pornography, which violated Section 2252(a)(4)(b). 
    \item Presentence report concluded that Lockhart was subject to Section 2252(b)(2)'s mandatory minimum because of his prior conviction related to aggravated sexual abuse, sexual abuse, or abusive sexual conduct involving a minor or ward. 
    \item Here, basically saying that minor/ward part of the statute only applies to the mention of abusive sexual conduct right before, which would construe the statute quite generally.
    \item Lockhart argued that minor/ward applied to all three of the listed crimes, and that his charges of sexual abuse involving an adult fell outside the scope. 
\end{itemize}

\rulebox{Section 2252(b)(2)}
{Defendants convicted of possessing child pornography in violation of 18 U.S.C. § 2252(a)(4) are subject to a 10–year mandatory minimum sentence and an increased maximum sentence if they have ``a prior conviction ... under the laws of any State relating to aggravated sexual abuse, sexual abuse, or abusive sexual conduct involving a minor or ward.''} 

\textbf{Question} Does the phrase ``involving a minor or ward'' modify all items in the list of predicate crimes, or only the item that precede it?

\textbf{Holding}
\begin{itemize}
    \item Held that the minor/ward clause only modifies the ``abusive sexual conduct'' phrase directly in front of it. 
    \item Therefore, Lockhart's previous indictment does apply in this case. 
\end{itemize}

\textbf{Discussion}

\begin{enumerate}
    \item \textbf{Rule of last antecedent vs Series-qualifier principle}
    \begin{itemize}
        \item \textbf{Majority:} the modifier applies only to ``abusive sexual conduct.''
        \item \textbf{Dissent:} the statute contains an integrated list, so the series-qualifier principle applies.
    \end{itemize}

    \item \textbf{Whether context overcomes the grammatical presumption}
    \begin{itemize}
        \item \textbf{Majority:} the list is structurally irregular, and Chapter 109A supports the last-antecedent reading.
        \item \textbf{Dissent:} the offenses are sufficiently parallel and related that ordinary English applies the modifier to all three.
    \end{itemize}

    \item \textbf{Statutory structure and Chapter 109A}
    \begin{itemize}
        \item \textbf{Majority:} Congress modeled the state predicates on Chapter 109A, which includes offenses involving adults and minors.
        \item \textbf{Majority:} Wouldn't make sense for there to be two different outcomes of almost identical state and federal crimes. 
        \item \textbf{Dissent:} the statutes do not precisely match; Chapter 109A contains four offenses, while Section 2252(b)(2) lists only three.
        \item \textbf{Dissent:}the majority placed too much weight on the similarities between the statutes, but the categories do not precisely mirror one another, and so cannot say Congress intended to reproduce Chapter 109A's structure in Section 2252(b)(2).
    \end{itemize}

    \item \textbf{Rule against surplusage}
    \begin{itemize}
        \item \textbf{Majority:} applying the modifier to all three offenses makes the list excessively redundant.
        \item \textbf{Dissent:} the majority's interpretation also creates redundancy and makes ``involving a minor or ward'' largely unnecessary.
    \end{itemize}

    \item \textbf{Rule of lenity}
    \begin{itemize}
        \item \textbf{Majority:} lenity applies only after ordinary interpretive tools fail; here, the text and structure provide a satisfactory interpretation.
        \item \textbf{Dissent:} because the statute remains reasonably ambiguous and imposes a criminal mandatory minimum, lenity should favor Lockhart.
    \end{itemize}
\end{enumerate}
}

\chapter{The Constitutional Position of Administrative Agencies (Regulation)}

\section{The Administrative State and the Delegation Question}

\section{Congressional Control of Agencies: The Legislative Veto}

\section{Appointment and Removal of Agency Officials} 

\section{Presidential Oversight of the Administration}

\chapter{The Administrative Rulemaking Process} 

\section{The Administrative Procedure Act and the Forms of Agency Action} 

\section{Notice \& Comment Rulemaking} 

\section{Alternatives to Notice \& Comment Rulemaking}

\subsection{Adjudication}

\subsection{Guidance Documents \& Interpretive Rules}

\subsection{Agreements}

\chapter{Judicial Review of Agency Action}

\section{Introduction to Judicial Review of Agency Decisions}

\section{Judicial (Non-)Regulation of Agency Procedures}

\section{“Hard Look” Review}

\chapter{Legal Interpretation in the Administrative State}

\section{Judicial Review of Agency Statutory Interpretation: Early Cases} 

\section{Judicial Review of Agency Statutory Interpretation: The Modern Framework}

\section{The Major Questions Doctrine}

\section{Judicial Review of Agency Statutory Interpretation: The New Framework}

\end{document}

